% from ega2-1-fr.tex
\oldpage[II]{199}
dans les homomorphismes (\hyperref[II.8.14.5.1-fr]{8.14.5.1}), on prend
$\mathcal{M}=\mathcal{S}$, on obtient l'homomorphisme de
$\mathcal{O}_Y$-Algèbres graduées
\begin{equation}
\label{II.8.14.5.3-fr}
  \alpha:\mathcal{S}\to\Gamma_*(\mathcal{S}_X)
\tag{8.14.5.3}
\end{equation}
déjà défini dans (\hyperref[II.3.3.2-fr]{3.3.2}), et qui fait donc de
$\Gamma_*(\mathcal{F})$ un $\mathcal{S}$-Module gradué~; les homomorphismes
(\hyperref[II.8.14.5.1-fr]{8.14.5.1}) définissent alors un homomorphisme (de
degré $0$) de $\mathcal{S}$-Modules gradués
\begin{equation}
\label{II.8.14.5.4-fr}
  \alpha:\mathcal{M}\to\Gamma_*(\operatorname{Proj}(\mathcal{M})).
\tag{8.14.5.4}
\end{equation}
\end{env}

\begin{env}[8.14.6]
\label{II.8.14.6-fr}
En général, pour un $\mathcal{S}_X$-Module gradué quasi-cohérent
$\mathcal{F}$, il n'est pas certain que le $\mathcal{S}$-Module gradué
$\Gamma_*(\mathcal{F})$ soit nécessairement quasi-cohérent. Considérons un
ouvert $X'$ de $X$ tel que la restriction $q':X'\to Y$ de $q$ à $X'$ soit
un morphisme \emph{quasi-compact}. Comme $q'$ est en outre séparé,
$q'_*(\mathcal{F}')$ est alors un $\mathcal{O}_Y$-Module quasi-cohérent pour
tout $\mathcal{O}_{X'}$-Module quasi-cohérent $\mathcal{F}'$
(\hyperref[I.9.2.2-fr]{I, 9.2.2, b)}). Nous poserons
\begin{equation}
\label{II.8.14.6.1-fr}
  \mathcal{S}_{X'}=\mathcal{S}_X|X'
  =\bigoplus_{n\in\mathbf{Z}}\mathcal{O}_X(n)|X'
\tag{8.14.6.1}
\end{equation}
et, pour tout $\mathcal{S}_{X'}$-Module gradué $\mathcal{F}'$,
\begin{equation}
\label{II.8.14.6.2-fr}
  \Gamma'_*(\mathcal{F}')=
  \bigoplus_{n\in\mathbf{Z}}q'_*(\mathcal{F}'_n).
\tag{8.14.6.2}
\end{equation}
La remarque précédente montre alors que si $\mathcal{F}'$ est un
$\mathcal{S}_{X'}$-Module quasi-cohérent, $\Gamma'_*(\mathcal{F}')$ est un
$\mathcal{S}$-Module gradué quasi-cohérent
(\hyperref[I.9.6.1-fr]{I, 9.6.1}).

On notera d'ailleurs que l'injection canonique $j:X'\to X$ est
\emph{quasi-compacte} puisque $q'=q\circ j$ l'est et que $q$ est séparé
(\hyperref[I.6.6.4-fr]{I, 6.6.4, (v)}). Par suite
$\mathcal{F}=j_*(\mathcal{F}')$ est un $\mathcal{S}_X$-Module gradué
quasi-cohérent pour tout $\mathcal{S}_{X'}$-Module gradué quasi-cohérent
$\mathcal{F}'$, et il résulte des définitions précédentes que l'on a
\begin{equation}
\label{II.8.14.6.3-fr}
  \Gamma'_*(\mathcal{F}')=\Gamma_*(\mathcal{F}).
\tag{8.14.6.3}
\end{equation}
Avec les mêmes hypothèses sur $X'$, pour tout $\mathcal{S}$-Module gradué
quasi-cohérent $\mathcal{M}$, nous poserons
\begin{equation}
\label{II.8.14.6.4-fr}
  \operatorname{Proj}'(\mathcal{M})
  =\operatorname{Proj}(\mathcal{M})|X'
\tag{8.14.6.4}
\end{equation}
qui est un $\mathcal{S}_{X'}$-Module gradué quasi-cohérent. L'homomorphisme
canonique
\[
  \operatorname{Proj}(\mathcal{M})\to
  j_*(\operatorname{Proj}'(\mathcal{M}))
\]
(\hyperref[0.4.4.3-fr]{0, 4.4.3}) donne donc un homomorphisme canonique
$\Gamma_*(\operatorname{Proj}(\mathcal{M}))\to
\Gamma'_*(\operatorname{Proj}'(\mathcal{M}))$ de $\mathcal{S}$-Modules
gradués, et par suite, en le composant avec
(\hyperref[II.8.14.5.4-fr]{8.14.5.4}), on obtient un homomorphisme canonique
fonctoriel (de degré $0$) de $\mathcal{S}$-Modules gradués quasi-cohérents
\begin{equation}
\label{II.8.14.6.5-fr}
  \alpha':\mathcal{M}\to
  \Gamma'_*(\operatorname{Proj}'(\mathcal{M})).
\tag{8.14.6.5}
\end{equation}
\end{env}

\begin{env}[8.14.7]
\label{II.8.14.7-fr}
Conservons sur $X'$ les hypothèses de
(\hyperref[II.8.14.6-fr]{8.14.6}), et soit $\mathcal{F}'$ un
$\mathcal{S}_{X'}$-Module gradué quasi-cohérent, de sorte que
$\operatorname{Proj}'(\Gamma'_*(\mathcal{F}'))$ est aussi un
$\mathcal{S}_{X'}$-Module gradué quasi-cohérent. Nous allons définir un
homomorphisme canonique fonctoriel (de degré $0$) de
$\mathcal{S}_{X'}$-Modules gradués
\begin{equation}
\label{II.8.14.7.1-fr}
  \beta':\operatorname{Proj}'(\Gamma'_*(\mathcal{F}'))\to\mathcal{F}'.
\tag{8.14.7.1}
\end{equation}

